% ID: SEL-D0039-OA-E01
% ESTADO: cotejada. Segunda pasada visual y borrador compilado revisados.
% Fuente: Selectividad/matematicas_aplicadas/Ua4MATEMÁTICAS APLICADAS examen.pdf
% SHA-256: 78ff6bdd48cbac71e6cc47d6ddf3dcb9492540ecb9cbe2032ebfec02e39fb66d
% Página: 1; opción A; ejercicio original 1.

\textbf{(2.5 puntos)} Con motivo de su inauguración, una heladería quiere repartir dos tipos de tarrinas de helados. El primer tipo de tarrina está compuesto por 100 g de helado de chocolate, 200 g de helado de straciatella y 1 barquillo. El segundo tipo llevará 150 g de helado de chocolate, 150 g de helado de straciatella y 2 barquillos. Sólo se dispone de 8 kg de helado de chocolate, 10 kg de helado de straciatella y 100 barquillos.

\textquestiondown{}Cuántas tarrinas de cada tipo se deben preparar para repartir el máximo número posible de tarrinas?
